\section{Annexe : SvelteKit pour les non-initiés}
\label{sec:sveltekit}

Cette annexe explique les concepts SvelteKit utilisés dans l'application, à destination de lecteurs familiers avec HTML/CSS/JS mais pas avec les frameworks modernes.

\subsection{Qu'est-ce qu'un composant Svelte~?}

Un fichier \texttt{.svelte} est un composant autonome qui regroupe HTML, JavaScript et CSS en un seul fichier~:

\begin{lstlisting}[language=svelte, title=Exemple : un composant minimal]
<script>
  let count = $state(0);
</script>

<button onclick={() => count++}>
  Compteur : {count}
</button>

<style>
  button { font-size: 1.2rem; }
</style>
\end{lstlisting}

\begin{itemize}
  \item \texttt{<script>} : logique JavaScript.
  \item Corps HTML : le template, avec des expressions \texttt{\{...\}} pour les valeurs dynamiques.
  \item \texttt{<style>} : CSS \emph{scopé} --- il ne s'applique qu'à ce composant, pas aux autres.
\end{itemize}

\subsection{Réactivité Svelte~5}

Svelte~5 utilise des \emph{runes} (fonctions spéciales préfixées par \texttt{\$}) pour la réactivité~:

\subsubsection{\texttt{\$state} --- État local réactif}

Déclare une variable dont les modifications sont automatiquement reflétées dans le DOM.

\begin{lstlisting}[language=javascript]
let count = $state(0);   // reactif
count++;                  // le DOM se met a jour automatiquement
\end{lstlisting}

\textbf{Dans l'application :} \texttt{HtmlContent} utilise \texttt{let container = \$state(null)} pour stocker la référence DOM.

\subsubsection{\texttt{\$derived} --- Valeur dérivée}

Calcule une valeur qui se met à jour automatiquement quand ses dépendances changent.

\begin{lstlisting}[language=javascript]
let count = $state(0);
let doubled = $derived(count * 2);  // toujours = count * 2
\end{lstlisting}

\textbf{Dans l'application :} \texttt{WordDetails} utilise \texttt{\$derived} pour calculer les détails du mot sélectionné à partir des stores.

La variante \texttt{\$derived.by(() => ...)} permet des calculs plus complexes (multi-lignes).

\subsubsection{\texttt{\$effect} --- Effet de bord réactif}

Exécute du code quand ses dépendances réactives changent.

\begin{lstlisting}[language=javascript]
$effect(() => {
  const enabled = $accentEnabled;  // dependance trackee
  if (container) applyAccents(container);
});
\end{lstlisting}

S'exécute :
\begin{enumerate}
  \item Au montage initial du composant.
  \item À chaque fois qu'une dépendance trackée change.
\end{enumerate}

\textbf{Dans l'application :} \texttt{HtmlContent} utilise \texttt{\$effect} pour ré-appliquer les accents et les handlers de hover quand l'état change.

\subsubsection{\texttt{\$props} --- Props du composant}

Déclare les propriétés reçues du composant parent.

\begin{lstlisting}[language=javascript]
// Dans le composant enfant :
let { html, disableHover = false } = $props();

// Dans le composant parent :
<HtmlContent html={wordInfo.base_html} disableHover={true} />
\end{lstlisting}

\subsection{Stores Svelte}

Un \emph{store} est un conteneur de valeur partagé entre composants, sans avoir besoin de passer des props à travers l'arbre.

\begin{lstlisting}[language=javascript, title=Creer un store]
import { writable } from 'svelte/store';
export const count = writable(0);
\end{lstlisting}

\begin{lstlisting}[language=javascript, title=Utiliser un store]
// Ecriture :
count.set(42);

// Lecture dans un composant (syntaxe $) :
<p>Valeur : {$count}</p>

// Lecture hors composant :
import { get } from 'svelte/store';
const val = get(count);
\end{lstlisting}

La syntaxe \texttt{\$count} (préfixe \texttt{\$}) est un raccourci Svelte qui s'abonne automatiquement au store et réagit aux changements. Hors d'un composant (dans une fonction utilitaire), il faut utiliser \texttt{get(store)}.

\textbf{Dans l'application :} Les stores sont définis dans \texttt{src/lib/stores/} et utilisés partout --- \texttt{\$wordData} dans les composants, \texttt{get(wordData)} dans \texttt{HtmlContent} (fonctions impératives).

\subsection{Routage basé sur le système de fichiers}

SvelteKit associe automatiquement les URLs aux fichiers dans \texttt{src/routes/}~:

\begin{center}
\begin{tabular}{ll}
  \toprule
  \textbf{Fichier} & \textbf{URL} \\
  \midrule
  \texttt{src/routes/+page.svelte} & \texttt{/} \\
  \texttt{src/routes/phrases/+page.svelte} & \texttt{/phrases} \\
  \bottomrule
\end{tabular}
\end{center}

Fichiers spéciaux dans chaque dossier de route~:
\begin{description}
  \item[\texttt{+page.svelte}] Le contenu de la page.
  \item[\texttt{+layout.svelte}] Template partagé par toutes les pages du dossier (et sous-dossiers). Affiche \texttt{\{@render children()\}} là où le contenu de la page doit apparaître.
  \item[\texttt{+layout.server.js}] Code exécuté côté serveur (ou au build) pour charger des données. La fonction \texttt{load()} retourne un objet accessible dans le layout via \texttt{\$props().data}.
\end{description}

\subsection{Adapter-static et prérendu}

Par défaut, SvelteKit génère une application serveur. L'option \texttt{adapter-static} transforme toutes les pages en HTML statique au moment du build~:

\begin{lstlisting}[language=javascript, title=svelte.config.js]
import adapter from '@sveltejs/adapter-static';

export default {
  kit: {
    adapter: adapter({
      pages: 'build',
      assets: 'build',
      strict: true
    })
  }
};
\end{lstlisting}

Avec \texttt{export const prerender = true} dans \texttt{+layout.server.js}, toutes les pages sont générées au build. Le résultat dans \texttt{build/} peut être déployé sur n'importe quel serveur statique (Apache, Nginx, GitHub Pages, etc.).

\subsection{Comparaison avant/après migration}

\begin{center}
\begin{tabular}{p{0.45\textwidth}p{0.45\textwidth}}
  \toprule
  \textbf{Vanilla JS (avant)} & \textbf{SvelteKit (après)} \\
  \midrule
  Scripts dans \texttt{scripts/*.js} avec manipulation DOM manuelle & Composants \texttt{.svelte} déclaratifs \\
  \midrule
  État global via variables globales & Stores Svelte (\texttt{writable}) \\
  \midrule
  \texttt{innerHTML} pour les tableaux & Templates Svelte avec \texttt{\{\#each\}}, \texttt{\{\#if\}} \\
  \midrule
  \texttt{document.querySelector} pour cibler les zones & Binding avec \texttt{bind:this} et refs \\
  \midrule
  Pages HTML séparées (\texttt{index.html}, \texttt{affiche\_phrase.html}) & Routes SvelteKit (\texttt{/}, \texttt{/phrases}) \\
  \midrule
  Chargement JSON via \texttt{fetch} au runtime & Prérendu : JSON intégré au build \\
  \midrule
  CSS global unique (\texttt{styles.css}) & CSS scopé par composant + \texttt{app.css} global \\
  \bottomrule
\end{tabular}
\end{center}
